\documentclass[11pt,a4paper]{article}
\usepackage[T1]{fontenc}
\usepackage[utf8]{inputenc}
\usepackage{lmodern,amsmath,amssymb}
\usepackage[margin=25mm]{geometry}
\usepackage[hidelinks]{hyperref}
\hypersetup{pdftitle={What a necessity argument for $h>0$ must supply: a unit and an indeterminacy},pdfauthor={navstokgap research working notes}}
\newcommand{\tightlist}{\setlength{\itemsep}{0pt}\setlength{\parskip}{0pt}}
\setlength{\emergencystretch}{3em}
\setcounter{secnumdepth}{0}
\begin{document}
\hypertarget{what-a-necessity-argument-for-h0-must-supply-a-unit-and-an-indeterminacy}{%
\section{\texorpdfstring{What a necessity argument for \(h>0\) must
supply: a unit and an
indeterminacy}{What a necessity argument for h\textgreater0 must supply: a unit and an indeterminacy}}\label{what-a-necessity-argument-for-h0-must-supply-a-unit-and-an-indeterminacy}}

\textbf{Result, 2026-09-27.} A positive universal floor on recorded
comparisons in Newton's Galileo comparison needs two separable
ingredients. \textbf{(U)} a universal action unit, and \textbf{(I)} an
indeterminacy of the record's back-action, since the floors of the
\href{planck-gap-paper.md}{Planck paper} (Theorems 5--6) bound
\emph{spreads} of disturbances, and a definite, later recoverable kick
has zero spread. The two have opposite status.

\begin{itemize}
\tightlist
\item
  \textbf{(U) follows from classical radiation thermodynamics} (Theorem
  U). Kirchhoff's universality, Wien's scaling law, finite equilibrium
  energy density and low-frequency equipartition force the equilibrium
  spectrum to contain a universal constant \(\beta\) with units of
  kelvin-seconds. With Boltzmann's constant it gives an action
  \(h=k_B\beta\), independent of every material constant. Its value is
  already in Wien's displacement constant \(b\): \(h=y_*k_Bb/c\), with
  \(y_*\) a pure number fixed by the spectral shape. No premise about
  the body, its mass or a commutator enters.
\item
  \textbf{(I) is impossible classically} (Theorem I). Any model with
  Liouville dynamics of body and pointers, product preparations obeying
  a prior density ceiling, and Bayesian conditioning on pointer readings
  admits records whose body posterior has arbitrarily small area. A
  premise that yields (I) must deny one of these three. Quantum
  kinematics denies the third; a universal unrecordable background
  (stochastic electrodynamics) denies the second; an epistemic
  restriction in Spekkens's sense imposes the ceiling directly. Newton's
  printed optics denies none, which is paper §8's \(\neg\)M3 as an
  instance.
\end{itemize}

So the necessity question becomes sharper than before. The existence of
a universal action unit is forced by nineteenth-century physics that
never mentions the body. What remains open is which of three explicit
classical structures a physically justified premise must give up for
that unit to act as a floor on records.

The candidate list and both theorems were proposed by a Fable reviewer
at the user's request (report in the session scratchpad, 2026-09-27),
after three external models (DeepSeek, Mimosa, Muse Spark 1.1) converged
on the Stoney unit \(k_e/c\), which the
\href{action-unit-dimensional-selection.md}{dimensional note} already
holds. The ingredients are classical and the theorems elementary; no
novelty is claimed for either. Their use here is to replace the
obstruction map of
\href{action-scale-obstructions.md}{action-scale-obstructions} by one
statement of what is missing.

\hypertarget{two-ingredients}{%
\subsection{1. Two ingredients}\label{two-ingredients}}

The floor of the Planck paper, Theorem 6, is

\[\frac s8\sum_j\Delta(\hat D_j)+\frac J2\sum_j\Delta(\hat X_j)
\ge\hbar\arcsin(1-2\epsilon),\]

with \(\Delta\) a spread of an undetermined disturbance. Its proof uses
Mandelstam--Tamm in the Bures angle and the Weyl relations, that is,
\(\hbar>0\) in two roles.

\begin{itemize}
\tightlist
\item
  \textbf{(U), the unit.} By Theorem A of the
  \href{action-unit-dimensional-selection.md}{dimensional note}, a floor
  shared by all masses and preparations is a pure number times a product
  of the model's fixed constants with action units. Classical mechanics,
  gravity and electrodynamics hold exactly one mass-independent such
  product, \(k_e/c=e^2/(4\pi\epsilon_0c)=\alpha\hbar\).
\item
  \textbf{(I), the indeterminacy.} The floor bounds spreads. A record
  whose kick is definite and later recoverable contributes nothing,
  whatever the size of the kick. For example, a charged probe passing a
  charged body at speed \(u<c\) and impact parameter \(b\) delivers the
  transverse impulse \(2n_1n_2k_e/(bu)\), so charge atomicity gives
  \(|\Delta p|\,b\ge2k_e/c=2\alpha\hbar\): a universal, mass-independent
  floor on the \emph{magnitude} of a record's kick. The kick is
  determined by the reading, so its spread is zero. The right unit gives
  the wrong kind of quantity.
\end{itemize}

\hypertarget{theorem-u-the-unit-from-radiation-thermodynamics}{%
\subsection{2. Theorem U: the unit from radiation
thermodynamics}\label{theorem-u-the-unit-from-radiation-thermodynamics}}

Let \(u(\nu,T)\,d\nu\) be the equilibrium energy density of
electromagnetic radiation between frequencies \(\nu\) and \(\nu+d\nu\)
at absolute temperature \(T\). Assume:

\begin{itemize}
\tightlist
\item
  \textbf{(K)} \(u\) is the same for every material enclosure (Kirchhoff
  1859, from the second law: two cavities at one temperature with
  different spectra would drive an engine through a colour filter);
\item
  \textbf{(W)} \(u(\nu,T)=\nu^3f(\nu/T)\) (Wien 1893, from adiabatic
  compression of a reflecting cavity and the second law
  (\href{https://doi.org/10.1007/978-3-663-13885-3_12}{Wien 1893,
  reprint}, metadata));
\item
  \textbf{(S)} \(\int_0^\infty u\,d\nu<\infty\) for every \(T>0\)
  (Stefan's measured law; Boltzmann's \(T^4\) theorem from Maxwell's
  radiation pressure);
\item
  \textbf{(R)} \(u(\nu,T)\to(8\pi\nu^2/c^3)\,k_BT\) as \(\nu/T\to0\)
  (equipartition of the low-frequency modes, the regime Rayleigh and
  Jeans describe).
\end{itemize}

\textbf{Theorem U.} Under (K), (W), (S), (R) there is a constant
\(\beta>0\) with units of kelvin-seconds, independent of every material
constant, such that
\(f(\nu/T)=(8\pi k_B/c^3)(T/\nu)\,\varphi(\beta\nu/T)\) with \(\varphi\)
a function of one dimensionless variable, \(\varphi(0)=1\). Hence
\(h:=k_B\beta\) is a universal action, and

\[\sigma'=\frac{8\pi k_B^4}{h^3c^3}\int_0^\infty y^2\varphi(y)\,dy,
\qquad h=\frac{y_*k_Bb}{c},\]

where \(\sigma'\) is the Stefan constant of the energy density, \(b\) is
Wien's displacement constant (\(\lambda_{\max}T=b\) for the peak of the
spectrum per unit wavelength), and \(y_*\) is the pure number at which
that peak sits.

\emph{Proof.} By (W), \(u\) depends on \(\nu\) and \(T\) through
\(\nu^3\) and the ratio \(\nu/T\), whose units are
\(1/({\rm s\cdot K})\). By (K) the only constants \(u\) may contain are
universal ones, and (R) displays two, \(c\) and \(k_B\). From \(c\),
\(k_B\) and \(\nu/T\) no dimensionless combination exists, because \(c\)
contains a length and \(k_B\) an energy, and neither appears in
\(\nu/T\). If \(u\) contained no further constant, \(f\) would be a pure
power, and (R) would fix it to \(f=(8\pi k_B/c^3)(T/\nu)\), whose
frequency integral diverges, contradicting (S). So \(u\) contains a
further universal constant. By (W) it can enter only through a
dimensionless multiple of \(\nu/T\), so it has units of kelvin-seconds;
call it \(\beta\). Then \(f\) has the displayed form, (R) gives
\(\varphi(0)=1\), and substituting \(y=\beta\nu/T\) in the integral
gives \(\sigma'\). The peak of \(u_\lambda=u\,c/\lambda^2\) lies at a
fixed value \(y_*\) of \(\beta c/(\lambda T)\), so
\(\lambda_{\max}T=\beta c/y_*\), which is the expression for
\(h=k_B\beta\). \(\square\)

With the spectral shape of Planck 1900
(\href{https://doi.org/10.1002/andp.19003060105}{Ann. Phys. 306, 69},
metadata), \(\varphi(y)=y/(e^y-1)\), the integral is \(\pi^4/15\),
\(y_*\approx4.965\), and \(h\) is Planck's constant. The theorem needs
no shape: any finite spectrum obeying (K), (W), (R) contains a universal
action, fixed up to a pure number by the measured \(b\), \(c\) and
\(k_B\). Planck formed his natural units from exactly this constant in
1899--1900, before the quantum hypothesis; the theorem states why a
constant had to be there.

Two remarks keep the statement honest. First, \(k_B\) converts kelvins
to energy, and its value needs atomism (the gas constant divided by
Avogadro's number); the universal constant of (K)--(W) alone is
\(\beta\). Second, with \(h\) and \(k_e/c\) both available, Theorem A of
the dimensional note no longer fixes a unique unit: any universal floor
is \(h\,G(\alpha)\) with \(\alpha=k_e/(hc/2\pi)\) dimensionless.
Uniqueness of the unit requires that the floor not depend on the charge,
which holds for records made with neutral probes.

\textbf{Bridge to the Opticks.} Newton's \emph{Opticks} holds, per
colour, an interval of fits \(\Lambda\) and a corpuscle whose
\(\Lambda p\) is invariant under refraction (paper §§7--8), but no
universality across colours: on his ordering of corpuscle sizes,
\(\Lambda p\) decreases toward the violet with no least colour. Theorem
U supplies exactly the missing cross-colour constant. Reading \(h\) as a
per-corpuscle impulse, \(p=h/\lambda\), needs more: Einstein's
fluctuation argument applied to an exponential spectral tail, which is
the quantization of the field. So the unit enters honestly through
(K)--(R), and quanta enter through the tail.

\hypertarget{theorem-i-no-classical-ceiling-on-posteriors}{%
\subsection{3. Theorem I: no classical ceiling on
posteriors}\label{theorem-i-no-classical-ceiling-on-posteriors}}

\textbf{Theorem I.} Consider a classical model with (i) Liouville
(Hamiltonian) dynamics of the body and its pointers, (ii) product
preparations, each factor with a prior phase-space density at most
\(1/X\), and (iii) records formed by Bayesian conditioning on a pointer
coordinate read to finite precision. Then for every \(X>0\) and
\(\eta>0\) there is a protocol whose body posterior satisfies
\(\Delta q\,\Delta p<\eta\), with \(\Delta\) the width of the support.

\emph{Proof.} Let the body prior be uniform on \([0,a]\times[0,X/a]\)
and the first pointer \((y,p_y)\) uniform on
\([0,\epsilon]\times[0,X/\epsilon]\); both have density \(1/X\). The
impulsive coupling with Hamiltonian \(g\,q\,p_y\) acting for unit time
maps \(y\mapsto y+gq\) and \(p\mapsto p-gp_y\), leaving \(q\) and
\(p_y\) unchanged, and preserves the joint density by Liouville. Reading
\(y\) to precision \(\epsilon'\) confines \(q\) to width
\((\epsilon+\epsilon')/g\). A second pointer \((z,p_z)\), uniform on
\([0,\epsilon_2]\times[0,X/\epsilon_2]\) and coupled by
\(g_2\,p_y\,p_z\), maps \(z\mapsto z+g_2p_y\) and kicks only \(y\),
which is already read. Reading \(z\) to precision \(\epsilon_2'\)
confines \(p_y\) to width \((\epsilon_2+\epsilon_2')/g_2\). The body's
momentum is \(p-gp_y\), so its posterior width is at most
\(X/a+g(\epsilon_2+\epsilon_2')/g_2\), and

\[\Delta q\,\Delta p\le\frac{X(\epsilon+\epsilon')}{ag}
+\frac{(\epsilon+\epsilon')(\epsilon_2+\epsilon_2')}{g_2},\]

which is below \(\eta\) for \(g_2\) large and \(a\) large, or
\(\epsilon\), \(\epsilon'\) small. The prior ceiling only widens the
unread conjugates \(p_y\) and \(p_z\). \(\square\)

The second pointer reads the first pointer's kick; this is paper §8's
countermodel (two position records recover the momentum) in mechanical
form, and each classical entry of the obstruction map is a corner of it.
The recoverable-delayed-record results C068--C069 are the corner
\(a\to\infty\), \(\epsilon'\to0\).

\textbf{What (I) must deny.} A premise that yields a universal posterior
ceiling, and with it paper Theorem 6 with the ceiling in place of
\(\hbar\), must deny (i), (ii) or (iii).

\begin{itemize}
\tightlist
\item
  \emph{Quantum kinematics} denies (iii): there is no joint phase-space
  density on which to condition, and every read is also a kick.
\item
  \emph{A universal background that no record can fully capture} denies
  (ii): every preparation is correlated with it. Stochastic
  electrodynamics is of this kind; Lorentz invariance fixes the spectrum
  of a random classical background as \(\propto\nu^3\) up to one
  constant with action units
  (\href{https://doi.org/10.1098/rspa.1963.0220}{Marshall 1963};
  \href{https://doi.org/10.1103/PhysRevD.11.790}{Boyer 1975}; metadata).
  The ceiling then needs the extra clause that records cannot read the
  background itself, since otherwise Theorem I applies to the enlarged
  system.
\item
  \emph{An epistemic restriction} imposes the ceiling on every state of
  knowledge. Bartlett, Rudolph and Spekkens show that Liouville
  mechanics with such a restriction reproduces Gaussian quantum
  mechanics (\href{https://doi.org/10.1103/PhysRevA.86.012103}{PRA 86,
  012103, 2012}, metadata;
  \href{https://doi.org/10.1103/PhysRevA.75.032110}{Spekkens 2007} for
  the discrete toy theory). This is the posterior ceiling as a
  principle, with its value free; Theorem U then fixes the value up to a
  number.
\item
  \emph{Newton's printed optics} denies none: its fits are determinate
  and its records conditionable, which is why paper §8 finds that what
  Newton printed entails \(\neg\)M3.
\end{itemize}

\hypertarget{consequence-for-state}{%
\subsection{4. Consequence for STATE}\label{consequence-for-state}}

STATE item 2 (Newton necessity) gains a sharp statement. The unit is
settled classically by Theorem U. The open task is a physically
justified premise that denies one of (i)--(iii) of Theorem I, of which
the three known routes are quantum kinematics, an unrecordable universal
background and an epistemic restriction. The historical test worth doing
next is the reviewer's second candidate: whether Newton's inflexion
experiments (\emph{Opticks} Book III, Obs. 1--11) give a bending law
\(\theta b\approx C\Lambda\) per colour. If they do, Newton's own optics
denies (iii) for corpuscle records, per colour, and paper §8 must be
reworded to say that \(\neg\)M3 follows only for inflexion-free records.
\end{document}
